\chapter{Research Methodology}
\label{chap:methodology}

\section{Research Design}
\label{sec:research_design}

This thesis follows a design science research (DSR) paradigm \citep{hevner2004design}, which is well-suited for the development and evaluation of IT artifacts in applied organizational contexts. DSR emphasizes the creation of purposeful artifacts — in this case, a forecasting and allocation recommendation framework — and their rigorous evaluation against real-world problems.

% TODO: justify DSR, describe iteration cycles

\section{Data Sources and Data Preparation}
\label{sec:data}

\subsection{Data Sources}
% TODO: describe what data is used (synthetic, anonymized, real?)
% - Historical project records (WBS, phase data, staffing records)
% - Skill taxonomies / HR master data
% - Actual vs planned resource usage

\subsection{Data Preprocessing}
% TODO: cleaning, feature engineering, temporal alignment

\subsection{Data Quality and Limitations}
% TODO

% ─────────────────────────────────────────────────────────────────────────────
\section{Modeling Approach}
\label{sec:modeling}

\subsection{Demand Forecasting Model}
% TODO: model selection rationale, feature set, target variable definition

\subsection{Bottleneck Detection Algorithm}
% TODO: threshold-based, anomaly detection, or constraint violation detection

\subsection{Allocation Optimization Model}
% TODO: problem formulation

Let the set of projects be $\mathcal{P} = \{p_1, p_2, \ldots, p_n\}$, the set of skill categories $\mathcal{S} = \{s_1, s_2, \ldots, s_m\}$, and the planning horizon $\mathcal{T} = \{t_1, t_2, \ldots, t_T\}$.

\noindent For each project $p \in \mathcal{P}$, skill $s \in \mathcal{S}$, and time period $t \in \mathcal{T}$, let:
\begin{itemize}
  \item $d_{p,s,t}$ — forecasted demand for skill $s$ in project $p$ at time $t$ (in FTE)
  \item $c_{s,t}$ — available capacity of skill $s$ at time $t$ (in FTE)
  \item $x_{p,s,t}$ — allocated resource of skill $s$ to project $p$ at time $t$ (decision variable)
\end{itemize}

\noindent The optimization objective minimizes the total unmet demand:
\begin{equation}
  \min \sum_{p \in \mathcal{P}} \sum_{s \in \mathcal{S}} \sum_{t \in \mathcal{T}} \max(0,\; d_{p,s,t} - x_{p,s,t})
  \label{eq:objective}
\end{equation}

\noindent Subject to:
\begin{align}
  \sum_{p \in \mathcal{P}} x_{p,s,t} &\leq c_{s,t} & \forall s \in \mathcal{S},\; t \in \mathcal{T} \label{eq:capacity} \\
  x_{p,s,t} &\geq 0 & \forall p, s, t \label{eq:nonnegativity}
\end{align}

% TODO: extend with priority weights, hard constraints per project

% ─────────────────────────────────────────────────────────────────────────────
\section{Evaluation Strategy}
\label{sec:evaluation_strategy}

\subsection{Forecasting Accuracy Metrics}
% TODO: MAE, RMSE, MAPE, skill-level breakdown

\subsection{Bottleneck Detection Performance}
% TODO: precision, recall, lead time of detection

\subsection{Allocation Quality Metrics}
% TODO: resource utilization rate, unmet demand rate, portfolio delay reduction

\subsection{Baseline Comparisons}
% TODO: naive/manual baseline, classical statistical model
